% =====================================================================
%  LỜI MỞ ĐẦU VÀ TỔNG QUAN BÀI TẬP
% =====================================================================
\chapter*{Lời mở đầu và tổng quan bài tập}
\addcontentsline{toc}{chapter}{Lời mở đầu và tổng quan bài tập}
\markboth{Lời mở đầu}{Lời mở đầu}

Các mô hình ngôn ngữ lớn (Large Language Models -- LLM) đã thay đổi cách máy tính xử lý ngôn ngữ tự nhiên: chỉ với một
câu lệnh, mô hình có thể tóm tắt văn bản, viết mã, trả lời câu hỏi chuyên môn. Nhưng một mô hình nền tảng thuần túy vẫn
làm việc theo kiểu \emph{hỏi -- đáp} (prompt-and-response). Nó chờ người dùng hỏi rồi mới trả lời, tri thức dừng lại ở
thời điểm huấn luyện, và nó không tự mở cơ sở dữ liệu, không tự gọi một API hay chạy một đoạn mã để thay đổi trạng thái
của thế giới số.

Bước chuyển đang diễn ra là đưa mô hình nền tảng vào bên trong một \textbf{hệ tác tử thông minh} (AI Agent, Agentic AI):
một hệ thống phần mềm tự đặt kế hoạch, tự chọn công cụ, ghi nhớ ngữ cảnh và hành động liên tục để đạt mục tiêu người dùng
giao. Bài tập Agent 01 yêu cầu tìm hiểu khái niệm này từ ba nguồn tài liệu công nghệ: \textbf{Google Cloud}
\cite{google_agents}, \textbf{IBM Think} \cite{ibm_agents} và \textbf{McKinsey} \cite{mckinsey_agents}. Đồng thời, bài tập
yêu cầu thực hành một dạng tác tử cổ điển là \textbf{tác tử học tăng cường} theo tài liệu \textbf{GeeksforGeeks}
\cite{gfg_rl}: lập trình trên Jupyter Notebook, chụp lại màn hình mã nguồn (copy screen code) và màn hình kết quả
(copy screen output).

\section*{Mục tiêu}
\begin{enumerate}
  \item Trình bày định nghĩa, ba trụ cột kỹ thuật và sự tiến hóa của AI Agent theo Google Cloud và McKinsey.
  \item Phân tích chu trình nhận thức, các mô hình suy luận (ReAct, ReWOO và các mô hình khác), các loại tác tử, rủi ro và
        nguyên tắc quản trị theo IBM Think.
  \item Trình bày nền tảng học tăng cường theo GeeksforGeeks: bài toán thăm dò -- khai thác, chính sách
        $\varepsilon$-greedy, phương trình Bellman và các trường phái thuật toán.
  \item Xây dựng tác tử Q-Learning trên Jupyter Notebook giải bài toán robot điều hướng mê cung, chụp màn hình mã và kết quả,
        rồi phân tích quá trình hội tụ của chính sách.
  \item Liên hệ học tăng cường với các tác tử LLM hiện đại (RLHF, suy luận dạng cây, hệ đa tác tử).
\end{enumerate}

\section*{Bố cục báo cáo}
Chương~1 và Chương~2 trình bày bản chất của AI Agent và lộ trình tiến hóa từ chatbot đến mạng lưới tác tử. Chương~3 đi vào
kiến trúc kỹ thuật và các cơ chế suy luận theo IBM Think. Chương~4 quy chiếu tác tử hiện đại về năm loại tác tử kinh điển
của Russell \& Norvig \cite{russell_norvig} và khung Quá trình Quyết định Markov. Chương~5 trình bày lý thuyết học tăng
cường. Chương~6 tổng hợp các nhóm ứng dụng doanh nghiệp cùng rủi ro và cách quản trị. Chương~7 và Chương~8 là phần thực
nghiệm: mô tả bài toán, mã nguồn, ảnh chụp màn hình notebook và phân tích kết quả. Chương~9 bàn về sự giao thoa giữa LLM
Agent và học tăng cường; Chương~10 tổng kết.

\begin{phathien}{Cách tạo số liệu trong báo cáo}
Notebook \texttt{agent01\_rl\_qlearning.ipynb} được sinh và chạy tự động từ mã nguồn trong gói \texttt{ag1}. Khi chạy xong,
notebook ghi kết quả ra tệp JSON. Mọi con số thực nghiệm trong văn bản và bảng được đọc từ tệp này khi dựng báo cáo; các
đoạn mã in trong báo cáo được trích nguyên văn từ chính mã đã chạy. Ảnh chụp màn hình là notebook đó mở trong JupyterLab.
\end{phathien}
